\section{A \textit{SystemUnderTest}}

    A \textit{SystemUnderTest}, vagy röviden \textit{Sut}, a teszt keretrendszer interakcióját valósítja meg a
    tesztelés alatt álló fordítóprogrammal. Ez a komponens lehetőséget ad fordítóprogram futtatását személyre szabni. A
    különböző \textit{Sut} implementációk egy bővítmény betöltő rendszeren keresztül válnak elérhetővé a teszt
    keretrendszer számára.

    \subsection{A \textit{Sut} függvények}

        A \textit{Sut} implementáció gyakorlatilag négy absztrakt függvény megvalósítását jelenti, amelyeket az alábbi
        \textbf{\ref{fig:sut-functions}. ábra} mutat. A függvények argumentumait a tömörség érdekében elhagytam.

        \begin{figure}[ht]
            \centering
            \begin{minted}{python}
class Sut(ABC):
    @abstractmethod
    async def setup(self, ...) -> SubprocessResult: ...

    @abstractmethod
    async def compile(self, ...) -> SubprocessResult: ...

    @abstractmethod
    async def link(self, ...) -> SubprocessResult: ...

    @abstractmethod
    async def run(self, ...) -> SubprocessResult: ...
        \end{minted}
            \caption{\label{fig:sut-functions} A \textit{Sut} absztrakt függvényei}
        \end{figure}

        A lépések a definíciós sorrendjében vannak végrehajtva.

        Az első \textit{setup} lépés lehetővé teszi a tesztfuttatás előtt extra lépéseket végrehajtani, például
        fájlokat létrehozni, környezeti változókat beállítani, vagy akár extra függvénykönyvtárakat is elérhetővé lehet
        tenni a fordítóprogram számára. Az ezt követő lépések opcionálisak, a tesztesetek típusától függenek.

        A második \textit{compile} lépésben történik a tesztesetek fordítása. Itt lehet személyre szabni a
        fordítóprogram számára átadott parancssori argumentumokat, amiket a \textit{ConfigManager} objektum
        előállított. Itt még történhet transzformáció a paramétereken, de egy kanonikus implementáció csak az
        argumentumok megfelelő sorrendjét állítja elő.

        A harmadik \textit{link} lépésben a szerkesztés, \textit{(idegen szóval linkelés)} történik. Itt ad lehetőség a
        rendszer a fordítóprogram szerkesztőjének parancssori argumentumait beállítani.

        A negyedik \textit{run} lépésben a fordítoprogram kimenetét lehet futtatni egy arra alkalmas szimulátor
        programban, vagy akár a végrehajtó gépen direktben.

        Fontos kiemelni, hogy bár ezek a függvények egymás után vannak végrehajtva, nincs megkötve hogy a tényleges
        implementációk milyen fordítóprogram invokációkat állítanak össze. Például ki is hagyhatnak lépéseket.

    \subsection{A \textit{CacheTarget} rendszer}

        A \textit{CacheTarget} rendszer lehetőséget ad a \textit{Sutok} \textit{run} függvényei számára, hogy
        eredményeiket egy gyorsítótárba helyezzék. Ez a gyorsítótár gyakorlatilag egy \textit{SQLite} adatbázis ahol
        egy kalkulált \textit{hashhez} hozzá van rendelve a futás eredménye.

        A rendszer olyan fordítóprogramokkal működik amelyek kimenetében \textit{ELF} formátumú tárgykódot generálnak.
        A rendszer megvizsgálja a kimeneti \textit{ELF} fájl fejlécét és szegmenseit és aggregálja a bájtokat. Ehhez az
        aggregációihoz még hozzáveszi a fordítoprogram bementi argumentumai és a \textit{STDIN} bemenetére érkező
        bájtokat, majd a végeredményből egy \textit{SHA-512 hasht}~\cite{sha512} generál.

        Az adatbázis ezt a \textit{hasht} és futáseredményt az alábbi \textbf{\ref{fig:cache-schema}. ábrán} látható
        táblában tárolja.

        \begin{figure}[ht]
            \centering
            \begin{tabular}{ll}
                \toprule
                \multicolumn{2}{c}{\textit{CacheTable}} \\
                \cmidrule{1-2}
                SQL type         & Name                 \\
                \midrule
                \textit{TEXT}    & test\_case           \\
                \textit{TEXT}    & hash                 \\
                \textit{TEXT}    & argv                 \\
                \textit{TEXT}    & stdin                \\
                \textit{INTEGER} & timeout              \\
                \textit{TEXT}    & stdout               \\
                \textit{INTEGER} & exit\_code           \\
                \textit{TEXT}    & stderr               \\
                \textit{INTEGER} & elapsed\_time        \\
                \bottomrule
            \end{tabular}
            \caption{\label{fig:cache-schema} A \textit{CacheTarget} adatbázis séma}
        \end{figure}

        A rendszer előnye, hogy a sokáig tartó futtatandó tesztesetek futásidejét drasztikusan tudja csökkenteni.
        Kifejezetten akkor a leghasznosabb ha a kimeneti \textit{ELF} futtatható szimulátorban fut, mert a szimulátorok
        természetükből adódóan lassabbak mint egy natív futtatás.

        A rendszer hátránya azonban az, hogy az elosztott futtatás során a kliensek nem osztják meg a gyorsítótáraikat
        egymással, ugyanis azok lemezen, \textit{SQLite} adatbázisokban vannak tárolva. Ha a futás során a szerver úgy
        osztja ki a teszteseteket, hogy egyik kliensen sem kerülne futtatásra ugyan az a teszt különböző
        tesztcsomagokból, akkor gyakorlatilag haszontalan arra a futásra.

        Ugyanakkor a kliensek megtarthatják a gyorsítótáraikat így jövőbeli futások alkalmával hasznosíthatóak.

    \subsection{A bővítmény rendszer}

        A \textit{Sut} komponensek dinamikusan kerülnek betöltésre. Ez lehetővé teszi, hogy a teszt keretrendszer
        csomagján kívülről is definiálható legyen egy \textit{Sut}. Ez a folyamat két lépésből áll, először
        regisztrálni kell a bővítményeket, majd be kell őket tölteni a \textit{Pythonban} szabványos \textit{importlib}
        rendszeren keresztül

        \subsubsection{Bővítmény regisztráció}

            Ennek megvalósítására a \textit{Python} nyelv által nyújtott \textit{metaclass} konstrukciót alkalmazta.

            A \textit{Python metclassok} egy fajta futásidejű formái a metaprogramozásnak, tehát futásidőben lehet a
            típusokat manipulálni. A bővítmény rendszer úgy használja fel ezt a konstrukciót, hogy amikor egy
            \textit{Sut} implementáció örököl a \textit{Sut} ősosztályból, akkor az ősosztály feljegyzi azt a típust
            egy asszociatív adatstruktúrába.

            Régebbi \textit{Python} verziókban szükségszerűen definiálni kellett egy osztályt ami a \textit{type}
            beépített típusból örökölt. Azonban sokkal elegánsabb és újabb megoldás az \\ \verb|__init_subclass__|
            függvény. A \textbf{\ref{fig:subclass-hook}. ábrán} látható ennek implementációja.

            \begin{figure}[ht]
                \centering
                \begin{minted}{python}
class Sut(ABC):
    def __init_subclass__(cls, target_name: str, **kwargs: Any) -> None:
        super().__init_subclass__(**kwargs)
        cls._SUBCLASSES[target_name] = cls  # storing the child class

    @classmethod
    def get_plugin(cls, target_name: str) -> type["Sut"] | None:
        return cls._SUBCLASSES.get(target_name)
            \end{minted}
                \caption{\label{fig:subclass-hook} A bővítmény regisztrálás implementációja}
            \end{figure}

            A \textit{Sut} implementációknak egyszerűen annyi dolguk van, hogy az öröklés során tovább adják a
            \verb|target_name| nevesített paramétert az ősosztálynak.

        \subsubsection{Bővítmény betöltés}

            A bővítmény betöltés a \textit{SutPluginLoader} osztályon keresztül történik. A bővítmény kezelő a bejövő
            modul mappájában egy \textit{ssplugin.py} nevű fájlt keres. Mivel a bővítmények szabványos \textit{Python}
            csomagok így nem elég csupán az elérési útvonalát ismerni a bővítmény fájlnak, hanem végre is kell hajtani
            azt.

            Erre többek között azért van szükség mert a \textit{Python} dinamikus így ameddig egy fájl nincs betöltve,
            addig a benne lévő definíciók nem is léteznek. Továbbá a globális szkópban lévő kódot is futtatni kell
            különben a modul használhatatlan lenne. A modul betöltését és végrehajtást a
            \textbf{\ref{fig:module-loder}. ábra} mutatja be.

            \begin{figure}[ht]
                \centering
                \begin{minted}{python}
class SutPluginLoader:
    def _load_module(self) -> None:
        # Error handling code omitted for brevity ...
        td: Final = self.target_dir
        spec: Final = spec_from_file_location(
            name=f"{td.parent.name}.{td.name}.ssplugin",
            location=td / "ssplugin.py",
        )
        module: Final = module_from_spec(spec)
        loader.exec_module(module)

    def load_sut_plugin(self) -> type[Sut]:
        self._load_module()
        if (plugin := Sut.get_plugin(self.target_dir.name)) is not None:
            return plugin
        # Error handling code omitted for brevity ...
            \end{minted}
                \caption{\label{fig:module-loder} A \textit{SutPluginLoader} implementációja.}
            \end{figure}

            A \verb|SutPluginLoader.load_sut_plgin| függvényével megkapjuk a dinamikusan betöltött \textit{Sut}
            implementációt.

            \pagebreak
    \subsection{Összegzés}

        Összességében a \textit{SystemUnderTest} egy nagyon funkciókban gazdag alrendszere a teszt keretrendszernek. A
        gyorsítótár implementációját leszámítva, az alkotó komponensei, koncepciójukat tekintve generikusak.

        A teszt keretrendszer nem tartalmaz olyan beépített bővítményt amely ne egy fordítóprogramot kezelne,
        ugyanakkor könnyen megvalósítható egy olyan \textit{Sut} bővítmény implementáció, amely nem csak
        fordítóprogramot tud tesztelni, hanem bármilyen parancssorból vezérelhető programot.

        Egy kézenfekvő példa egy \LaTeX{} fordítóprogramot tesztelő \textit{Sut} bővítmény.

        \begin{enumerate}
            \item A bővítmény a \textit{setup} lépésben beállítja megfelelő munka könyvtárat ugyanis a dokumentum
                  fordítók erre érzékenyek.
            \item A \textit{compile} lépésben végrehajtja az első futást ahol a referenciák még üresek.
            \item A \textit{link} lépésben végrehajtja a második futást ahol elkészül a teljes dokumentum \textit{DVI}
                  vagy \textit{PDF} formátumban.
            \item A \textit{run} lépésben pedig egy nyelvtani korrektor program ellenőrzi a generált dokumentumot.
        \end{enumerate}
